% F5 -- Where a day is cut, and how a single position is scored.
\begin{figure}[!htbp]
\centering
\fitwidth{%
\begin{tikzpicture}[font=\small,
  ar/.style={-{Stealth[length=4pt]},clrule},
  ctx/.style={draw=clrule,fill=clsoft,rounded corners=2pt,inner sep=4pt,
              font=\scriptsize,text width=3.1cm,align=center},
  prd/.style={draw=clmodel,fill=clmodel!8,rounded corners=2pt,inner sep=4pt,
              font=\scriptsize,text width=2.2cm,align=center},
]
% ---- (a) the split ------------------------------------------------------
\node[anchor=west,font=\footnotesize\bfseries,text=clmodel] at (0,1.05)
  {(a) Each person's day is cut in two, once and for all};
\fill[clmodel!22] (0,0) rectangle (11.1,0.65);
\fill[clink!25]  (11.1,0) rectangle (13.9,0.65);
\draw[clrule] (0,0) rectangle (13.9,0.65);
\node[font=\footnotesize] at (5.55,0.33)
  {\textbf{THE PAST}, 80\,\% of the day, given to the model};
\node[font=\footnotesize] at (12.5,0.33) {\textbf{THE REST}, 20\,\%};
\node[font=\scriptsize,anchor=north,text=clrule] at (6.95,-0.10)
  {where this cut falls is called the \emph{context fraction}; several
   experiments move it from 30\,\% to 90\,\%};

% ---- (b) scoring one position at a time ---------------------------------
\node[anchor=west,font=\footnotesize\bfseries,text=clmodel] at (0,-1.20)
  {(b) The rest is then answered one moment at a time, never in one go};
\begin{scope}[yshift=-1.60cm]
  \node[ctx] (b0) at (2.6,-0.35)  {given events 1 to 40};
  \node[prd] (p0) at (5.9,-0.35)  {name event 41};
  \node[ctx] (b1) at (2.6,-1.10)  {given events 1 to 41};
  \node[prd] (p1) at (5.9,-1.10)  {name event 42};
  \node[ctx] (b2) at (2.6,-1.85)  {given events 1 to 42};
  \node[prd] (p2) at (5.9,-1.85)  {name event 43};
  \draw[ar] (b0) -- (p0); \draw[ar] (b1) -- (p1); \draw[ar] (b2) -- (p2);
  \node[font=\scriptsize,anchor=west,text width=6.6cm,align=left] at (7.6,-1.10)
    {Each time, the model is handed the person's \emph{real} past, not its own
     previous answers. So it can neither see the future nor go astray on its own
     earlier mistakes. The one-line rule answers exactly the same moments, and
     simply repeats the last antenna in the history it was given.};
\end{scope}
\end{tikzpicture}}
\caption{Where the day is cut, and how each moment is answered}
\label{fig:protocol}
\figtext{%
\figpara{Panel (a) is the cut; panel (b) is how each moment is answered.}{Each
test person's day is split once into a part the model may read and a part it has
to produce, the default being 80/20. Where that cut falls is part of the
\emph{measurement} rather than a property of the model, which is why it is drawn
here (Chapter~\ref{ch:first}). The model is then not asked to invent a whole
afternoon: it is asked for one antenna, told what really happened, and asked for
the next, so it can neither see the future nor compound its own errors. The
one-line rule is run through exactly the same moments with exactly the same
history and simply answers the last antenna it was shown, which is what makes the
two numbers comparable at all.}}
\end{figure}
